%\documentclass[final, onecolumn, 3p]{elsarticle}
\documentclass[..]{../IEEEphot}
\usepackage{array,multirow}

%******************************************
%*************** My packages **************
%******************************************
\usepackage{amsmath}
\usepackage{amssymb}
\usepackage{booktabs}
\usepackage[utf8]{inputenc}
\usepackage{longtable}
\usepackage{times}
\usepackage{textcomp}
\usepackage[svgnames]{xcolor}
\usepackage{scalerel}
\usepackage{keystroke}
%******************************************
%******************************************
%******************************************


%******************************************
%*************** My commands **************
%******************************************

%Circled characters
\makeatletter
\usepackage{tikz}
\usetikzlibrary{calc}
\newcommand*\circled[1]{\tikz[baseline=(char.base)]{
    \node[shape=circle, draw, inner sep=1pt, 
        minimum height={\f@size*1.6},] (char) {\vphantom{WAH1g}#1};}}
\makeatother
%******************************************
%******************************************
%******************************************

\newcommand{\param}[1]{\colorbox{LightSlateGray}{\color{Navy}\texttt{\textbf{#1}}}}

\newcolumntype{z}[1]{>{\centering\arraybackslash}m{#1}}% a centered m-type column
\renewcommand*\multirowsetup{\centering}% default is \raggedright

\def\msquare{\mathord{\scalerel*{\Box}{gX}}}



\begin{document}

\title{AutoCAD$^{\text{\textregistered}}$ cheat sheet}

\author{Francesco~Bianconi}

\affil{Department of Engineering\\ Università degli Studi di Perugia \\ Via Goffredo Duranti, 93 -- 06125 Perugia (Italy)\\ \texttt{bianco@ieee.org} \\} 

\maketitle
% * <john.hammersley@gmail.com> 2015-02-09T12:07:31.197Z:
%
%  Click the title above to edit the author information and abstract
%
\thispagestyle{empty}

\vspace{-1.0cm}
\noindent Up-to-date with AutoCAD 2023 (some commands may not work with previous versions). Last revision: \today.

\tableofcontents
	
\clearpage

\section{Reference systems and coordinate insertion}
\begin{longtable}{m{.3\linewidth}>{\centering\arraybackslash}p{.3\linewidth}>{\centering\arraybackslash}p{.3\linewidth}}
\toprule
\multirow{2}{*}{Reference system} & \multicolumn{2}{c}{Insertion mode} \\
\cmidrule{2-3}
 & Cartesian & Polar \\
\midrule
Global (WCS) & \multicolumn{1}{c}{*$x$,$y$} & \multicolumn{1}{c}{*$\rho<\theta$} \\
User (UCS) & \multicolumn{1}{c}{$x$,$y$} & \multicolumn{1}{c}{$\rho<\theta$} \\
Local & \multicolumn{1}{c}{@$x$,$y$} & \multicolumn{1}{c}{@$\rho<\theta$} \\
\bottomrule
\end{longtable}

\section{Creating graphical objects}

\begin{center}
\begin{longtable}{m{.1\linewidth}m{.2\linewidth}m{.3\linewidth}m{.3\linewidth}}
\toprule
    \multicolumn{2}{c}{\bfseries Command} &
    \multicolumn{1}{c}{\bfseries Effect} &
    \multicolumn{1}{c}{\bfseries Main parameters} \\
\midrule
\texttt{\_arc} & \includegraphics[width = 0.8\linewidth, keepaspectratio]{../images/jpg/_arc.jpg}
& Creates a circular arc. & 
\param{C}: Accepts the centre as the first input (instead of the starting point).
\\			
\midrule
\texttt{\_circle} & \includegraphics[width = 0.8\linewidth, keepaspectratio]{../images/jpg/_circle.jpg} 
& Creates a full circle. & 
\param{D}: Accepts the diameter instead of the radius.

\param{3P}: Generates a circle based on three points.

\param{2P}: Generates a circle based on two diametrically opposite points.

\param{T}: Generates a circle tangent to three objects.
\\	

\midrule
\texttt{\_ellipse} & \includegraphics[width = 0.8\linewidth, keepaspectratio]{../images/jpg/_ellipse.jpg}
& Creates a full ellipse or an elliptical arc. & 
\param{A}: Generates an elliptical arc (instead of a full ellipse).

\param{C}: Generates an ellipse or an elliptical arc starting from the centre.
\\
\midrule
\texttt{\_hatch} & \includegraphics[width = 0.8\linewidth, keepaspectratio]{../images/jpg/_hatch.jpg} & Hatches a closed region & 
\\
\midrule
\texttt{\_line} & \includegraphics[width = 0.8\linewidth, keepaspectratio]{../images/jpg/_line.jpg} & Generates a single line or a polyline. & 
\param{C}: Closes the polyline (only available when there the polyline has at least three points). \\
\midrule
\texttt{\_point} & \includegraphics[width = 0.8\linewidth, keepaspectratio]{../images/jpg/_point.jpg} & Generates one point. &  \\	
\midrule
\texttt{\_polygon} & \includegraphics[width = 0.8\linewidth, keepaspectratio]{../images/jpg/_polygon.jpg} & Generates an inscribed or circumscribed regular polygon. & 
\param{C}: Generates a circumscribed polygon.

\param{I}: Generates an inscribed polygon.  \\
\midrule
\texttt{\_rectang} & \includegraphics[width = 0.8\linewidth, keepaspectratio]{../images/jpg/_rectang.jpg} & Generates a rectangle. & 
\param{C}: Defines the dimensions of the chamfers (for a rectangle with chamfered vertices).

\param{F}: Defines the radius of the fillets (for a rectangle with rounded vertices).

\param{D}: Defines the width and height of the rectangle (available once the first corner is specified).

\\
\bottomrule
\end{longtable}
\end{center}

\clearpage

\section{Editing graphical objects}

\begin{center}
\begin{longtable}{m{.2\linewidth}m{.2\linewidth}m{.25\linewidth}m{.25\linewidth}}
\toprule
    \multicolumn{2}{c}{\bfseries Command} &
    \multicolumn{1}{c}{\bfseries Effect} &
    \multicolumn{1}{c}{\bfseries Main parameters} \\
\midrule
\texttt{\_arrayclassic} & \includegraphics[width = 0.8\linewidth, keepaspectratio]{../images/jpg/_arrayclassic.jpg} & Creates copies of objects along a linear or circular array. & 
\\
\midrule
\texttt{\_chamfer} & \includegraphics[width = 0.8\linewidth, keepaspectratio]{../images/jpg/_chamfer.jpg}  
& Chamfers two rectilinear edges. The two edges (or their extensions) must intersect. & 
\param{D}: Sets the chamfer distance from the intersecting point of the first and second object. The input order reflects the order whereby the entities are selected.
\\
\midrule
\texttt{\_copy} & \includegraphics[width = 0.8\linewidth, keepaspectratio]{../images/jpg/_copy.jpg} & Creates copies of objects. & 
\param{D}: Sets the displacement vector
\\	
\midrule
\texttt{\_erase} & & Deletes the selected objects (same effect as pressing \keystroke{Canc}). & 
\\	
\midrule
\texttt{\_explode} & & Splits an object into separate parts (opposite effect than \texttt{\_join}). & 
\\	
\midrule
\texttt{\_extend} & \includegraphics[width = 0.8\linewidth, keepaspectratio]{../images/jpg/_extend.jpg} & Extends objects to meet the edges of other objects.  & 
\\	
\midrule
\texttt{\_fillet} & \includegraphics[width = 0.8\linewidth, keepaspectratio]{../images/jpg/_fillet.jpg} & Fillets the edges of two 2D objects by creating an arc of given radius. The two edges (or their extensions) must intersect. & 
\param{R}: Sets the fillet radius.
\\
\midrule
\texttt{\_intersect} & \includegraphics[width = 0.8\linewidth, keepaspectratio]{../images/jpg/_intersect.jpg} & Returns the intersection (Boolean \texttt{AND}) between two closed objects. The objects must be previously converted into regions through \texttt{\_region}. & 
\\
\midrule
\texttt{\_join} & & Joins lines, polylines, splines and adjacent arcs (opposite effect than \texttt{\_explode}). & 
\\
\midrule
\texttt{\_mirror} & \includegraphics[width = 0.8\linewidth, keepaspectratio]{../images/jpg/_mirror.jpg} & Mirrors objects across a given axis. & 
\param{N}: Keeps the source objects (creates a copy).

\param{Y}: Removes the source objects.
\\		
\midrule
\texttt{\_move} & \includegraphics[width = 0.8\linewidth, keepaspectratio]{../images/jpg/_move.jpg} & Moves objects. & 
\param{D}: Sets the displacement vector.
\\			
\midrule
\texttt{\_offset} & \includegraphics[width = 0.8\linewidth, keepaspectratio]{../images/jpg/_offset.jpg} & Offsets an object to a given distance.  & 
\\			
\midrule
\texttt{\_rotate} & \includegraphics[width = 0.8\linewidth, keepaspectratio]{../images/jpg/_rotate.jpg} & Rotates objects. & 
\param{C}: Keeps the original objects (creates a copy).
\\	
\midrule
\texttt{\_scale} & \includegraphics[width = 0.8\linewidth, keepaspectratio]{../images/jpg/_scale.jpg} & Applies a scale transform. & 
\param{C}: Keeps the original objects (creates a copy).
\\	
\midrule
\texttt{\_subtract} & \includegraphics[width = 0.8\linewidth, keepaspectratio]{../images/jpg/_subtract.jpg} & Returns the difference (Boolean \texttt{XOR}) between two closed objects. The objects must be previously converted into regions through \texttt{\_region}. & 
\\
\midrule
\texttt{\_trim} & \includegraphics[width = 0.8\linewidth, keepaspectratio]{../images/jpg/_trim.jpg} & Trims objects to meet the edges of other objects.  & 
\\	
\midrule
\texttt{\_union} & \includegraphics[width = 0.8\linewidth, keepaspectratio]{../images/jpg/_union.jpg} & Returns the union (Boolean \texttt{OR}) between two closed objects. The objects must be previously converted into regions through \texttt{\_region}. & 
\\
\bottomrule
\end{longtable}
\end{center}

\clearpage

\section{Inserting dimensions and managing their style}

\begin{center}
\begin{longtable}{m{.2\linewidth}m{.2\linewidth}m{.25\linewidth}m{.25\linewidth}}
\toprule
    \multicolumn{2}{c}{\bfseries Command} &
    \multicolumn{1}{c}{\bfseries Effect} &
    \multicolumn{1}{c}{\bfseries Main parameters} \\
\midrule
\texttt{\_dimlinear} & \includegraphics[width = 0.8\linewidth, keepaspectratio]{../images/jpg/_dimlinear.jpg} & Creates a linear dimension parallel to one of the axes of the UCS. & 
\\
\midrule
\texttt{\_dimaligned} & \includegraphics[width = 0.8\linewidth, keepaspectratio]{../images/jpg/_dimaligned.jpg} & Creates a linear dimension that is aligned with the origin points of the extension lines. & 
\\
\midrule
\texttt{\_dimdiameter} & \includegraphics[width = 0.8\linewidth, keepaspectratio]{../images/jpg/_dimdiameter.jpg} & Creates a diameter dimension with a diameter symbol ($\phi$) in front of it. & 
\\
\midrule
\texttt{\_dimradius} & \includegraphics[width = 0.8\linewidth, keepaspectratio]{../images/jpg/_dimradius.jpg} & Creates a radius dimension with a radius symbol (\texttt{R}) in front of it. & 
\\
\midrule
\texttt{\_dimstyle} & & Opens the window that allows to manage the dimension style. & 
\\
\bottomrule
\end{longtable}
\end{center}

\clearpage

\section{Inserting text and annotations}

\begin{center}
\begin{longtable}{m{.2\linewidth}m{.2\linewidth}m{.25\linewidth}m{.25\linewidth}}
\toprule
    \multicolumn{2}{c}{\bfseries Command} &
    \multicolumn{1}{c}{\bfseries Effect} &
    \multicolumn{1}{c}{\bfseries Principali parametri} \\
\midrule
\texttt{\_mtext} & \includegraphics[width = 0.8\linewidth, keepaspectratio]{../images/jpg/_mtext.jpg} & Inserts multi-line text. & 
\param{H}: Specifies the text height.

\param{J}: Sets text justification and flow in relation to the text boundary (see \texttt{\_text}).

\param{W}: Specifies the width of the text boundary.
\\
\midrule
\texttt{\_qleader} & \includegraphics[width = 0.8\linewidth, keepaspectratio]{../images/jpg/_qleader.jpg} & Creates a leader and leader annotation. \\
\midrule
\texttt{\_text} & \includegraphics[width = 0.8\linewidth, keepaspectratio]{../images/jpg/_text.jpg} & Inserts single-line text. & 
\param{J}: Sets text justification and flow in relation to the text boundary. Possible options:
\begin{itemize}
\item \param{TL} (Top left)
\item \param{TC} (Top centre)
\item \param{TR} (Top right)
\item \param{ML} (Middle left)
\item \param{MC} (Middle centre)
\item \param{MR} (Middle right)
\item \param{BL} (Bottom left)
\item \param{BC} (Bottom centre)
\item \param{BR} (Bottom right)
\end{itemize}
\item \param{W} Specifies the width of the text boundary; words wrap to fit the width.
\\
\midrule
\texttt{\_tolerance} & \includegraphics[width = 0.8\linewidth, keepaspectratio]{../images/jpg/_tolerance.jpg} & Creates symbols for geometric tolerances and datums. \\
\midrule
\texttt{\_centerline} & \includegraphics[width = 0.8\linewidth, keepaspectratio]{../images/jpg/_centerline.jpg} & Draws the axis between two parallel or non-parallel lines. & 
\\
\midrule
\texttt{\_centermark} & \includegraphics[width = 0.8\linewidth, keepaspectratio]{../images/jpg/_centermark.jpg} & Draws the centre mark of a circle or arc. & 
\\
\bottomrule
\end{longtable}
\end{center}

\subsection{Inserting special characters}

\vspace{\baselineskip}

\begin{center}
\begin{tabular}{lccc}
\toprule
\textbf{Property name} & \textbf{Symbol} & \textbf{AutoCAD sequence} & \textbf{Unicode sequence}\\
\midrule
Diameter & \o & \%\%c & \textbackslash U+2205 \\
\midrule
Envelope requirement & \circled{E} & -- & \textbackslash U+24BA \\
\midrule
Degrees & $^{\circ}$ & \%\%d & \textbackslash U+00B0 \\
\midrule
Plus/minus & $^{\pm}$ & \%\%p & \textbackslash U+00B1  \\
\midrule
Square & $\msquare$ & -- & \textbackslash U+25A1 \\
\bottomrule
\end{tabular}
\end{center}

\noindent
 

\subsection{Inserting special characters via \texttt{gdt} font}

\vspace{\baselineskip}

\begin{center}
\begin{tabular}{lcc}
\toprule
\textbf{Property name} & \textbf{Symbol} & \textbf{AutoCAD sequence} \\
\midrule
Diameter & \o & n \\
\midrule
Square & $\msquare$ & o \\
\midrule
\raisebox{0pt}{Counterbore} & \includegraphics[width = 10pt, keepaspectratio]{../images/jpg/hole-counterbore.jpg} & \raisebox{0pt}{v} \\
\midrule
\raisebox{0pt}{Countersink} & \includegraphics[width = 10pt, keepaspectratio]{../images/jpg/hole-countersink.jpg} & \raisebox{0pt}{w} \\
\midrule
\raisebox{5pt}{Deep/Depth} & \includegraphics[width = 10pt, keepaspectratio]{../images/jpg/hole-depth.jpg} & \raisebox{5pt}{x} \\
\midrule
Taper & \includegraphics[height=\fontcharht\font`\B, keepaspectratio]{../images/jpg/GDandT_Conical_Taper_Symbol.jpg} & y \\
\midrule
Maximum material requirement & \circled{M} & m \\
\midrule
Minimum material requirement & \circled{L} & l \\
\bottomrule
\end{tabular}
\end{center}

\clearpage

\section{Managing user interface, system variables and object properties}
\begin{center}
\begin{longtable}{m{.2\linewidth}m{.33\linewidth}m{.33\linewidth}}
\toprule
    \multicolumn{1}{c}{\bfseries Command} &
    \multicolumn{1}{c}{\bfseries Effect} &
    \multicolumn{1}{c}{\bfseries Main parameters} \\
\midrule
\texttt{\_dynmode} & Sets the dynamic input mode. & 
\param{1}: Turns dynamic input on.

\param{0}: Turns dynamic input off.
\\
\midrule
\texttt{\_layer} & Opens the window for creating, deleting and modifying layers. & 
\\
\midrule
\texttt{\_matchprop} & Applies the properties of a selected object to other objects (similar to paste format in Microsoft Word). & 
\\
\midrule
\texttt{\_navbardisplay} & Controls the display of the navigation bar. & 
\param{1}: Displays the navigation bar.

\param{0}: Hides the navigation bar.
\\
\midrule
\texttt{\_orthomode} & Constrains cursor movement parallel to UCS axes. & 
\param{1}: Turns orthogonal mode on.

\param{0}: Turns orthogonal mode off.
\\
\midrule
\texttt{\_osnap} & Opens the window for managing object snaps. & \includegraphics[width = 0.8\linewidth, keepaspectratio]{../images/png/_osnap_en.png}
\\
\midrule
\texttt{\_polarang} & Sets the polar angle increment for tracking. & 
\param{0}: No polar tracking.

\param{]0,90]}: Polar tracking by steps of the given value.
\\
\midrule
\texttt{\_ptype} & Opens the window for setting the display style and size of point objects.  & \includegraphics[width = 0.8\linewidth, keepaspectratio]{../images/png/_ptype_en.png} 
\\
\midrule
\texttt{\_qselect} & Selects objects based on filtering criteria. & \includegraphics[width = 0.8\linewidth, keepaspectratio]{../images/png/_qselect_en.png} 
\\
\midrule
\texttt{\_regen} & Regenerates the drawing in the current viewport. Recomputes the locations and visibility of all objects in the current viewport. & 
\\
\midrule
\texttt{\_redraw} & Refreshes the display in the current viewport. & 
\\
\midrule
\texttt{\_savetime} & Sets the automatic save time interval (in minutes). & 
\\
\midrule
\texttt{\_ucs} & Sets the origin and orientation of the UCS. & 
\\
\midrule
\texttt{\_ucsicon} & Controls the appearance of the UCS icon. & 
\param{ON}: Shows UCS icon.

\param{OF}: Hides UCS icon.

\param{OR}: Displays the icon at the origin (0,0,0) of the UCS.
\\
\midrule
\texttt{\_units} & Opens the dialogue box to control format, precision and other settings used to display coordinates, distances and angles.
\\
\bottomrule
\end{longtable}
\end{center}

\clearpage

\section{Keyboard shortcuts}

\begin{center}
\begin{longtable}{m{.2\textwidth}m{.7\textwidth}m{.1\textwidth}}
\toprule
    \multicolumn{1}{c}{\bfseries Key} &
    \multicolumn{1}{c}{\bfseries Effect(s)}\\
\midrule
\keystroke{Canc} & Deletes the selected objects. \\
\midrule
\keystroke{Ctrl}+\keystroke{B} & Turns the grid snap on/off. \\
\midrule
\keystroke{Ctrl}+\keystroke{F4} & Closes the current document. \\
\midrule
\keystroke{Ctrl}+\keystroke{O} & Opens an existing document. \\
\midrule
\keystroke{Ctrl}+\keystroke{Q} & Closes the application. \\
\midrule
\keystroke{Ctrl}+\keystroke{Z} & Undoes the previous action. \\
\midrule
\keystroke{Ctrl}+\keystroke{Y} & Redoes the previous action (reverses a previous undo). \\
\midrule
\keystroke{Ctrl}+\keystroke{9} & Hides/displays the command line. \\
\midrule
\keystroke{Esc} & 
Exits a command. 

Deselects all objects.
\\
\midrule
\keystroke{Invio} & 
Finalises a command, the input of a parameter or the selection of objects within one command. 

Repeats the latest command.
\\
\midrule
\keystroke{F3} & 
Turns object snap on/off.\\
\midrule
\keystroke{F8} & 
Turns orthogonal mode on/off.\\
\midrule
\keystroke{F10} & 
Turns polar tracking on/off.\\
\midrule
\keystroke{F11} & 
Turns object snap tracking on/off.\\
\midrule
\keystroke{F12} & 
Turns dynamic mode on (but not off).\\
\midrule
\keystroke{$\Uparrow$}, \keystroke{Shift} &
Removes objects from the current selection when pressed or held down while selecting objects. \\
\bottomrule
\end{longtable}
\end{center}

\end{document}